\section{Conclusion}
\label{sec:conclusion}

For an incremental, masked MPLS traffic-engineering problem with persistent
routing state, a masked contextual bandit beats MaskablePPO on every training
root and matches a per-interval MILP optimizer. The cause is the planning
horizon, not the learning algorithm: PPO with $\gamma=0$ matches the bandit,
while PPO with $\gamma=0.995$ settles into costly oscillations, and a
bootstrapped value learner does not improve on the bandit. A sequentiality
diagnostic shows why: almost all of the problem's non-myopic value lies in
anticipating exogenous change that the controller cannot observe, while the
predictable part is captured by the first interval's reward. Sequential RL
should therefore be justified by measured temporal coupling---delayed effects
or predictable dynamics---rather than assumed from the fact that the problem
is formally an MDP.
